A growing stream of blink and ocular-artifact research replaces hand-tuned thresholds with data-driven models. Convolutional neural networks have been used both to detect eye-blink artifacts and to remove them from EEG~\citep{jurczak2022implementation,gawande2022deep}, while transformer- and video-based architectures have pushed blink-detection accuracy further~\citep{fodor2023blinklinmult,hong2024robust}. Autoencoder and segmentation-denoising networks also target single-channel artifact removal~\citep{faisal2025eye,li2023segmentation}, and recent reviews catalogue the broader shift from conventional thresholding toward deep learning for physiological-artifact handling, including in wearable EEG~\citep{akshathraj2026physiological,arpaia2025systematic}. In driving and fatigue monitoring, machine-learning and CNN models detect drowsiness from blink and eye-state cues~\citep{cabot2024drowsiness,makhmudov2024real}. These data-driven detectors can be accurate, but they typically require labelled training data, more computation, and offer less transparency than a single interpretable threshold. This contrast motivates studying the lightweight thresholding stage itself, which frequently precedes or complements such pipelines, and in particular the epoch-structure estimation bias that this paper addresses.
